% !TeX root = main_long.tex
\section{RQ2: cross-implementation consistency and round-tripping}
\subsection{Method}
\label{sec:m-rq2}

RQ2 examines whether VCF Core can support a second conversion strategy and whether structured RDF retains enough information to recover the source observations.
We re-execute all RQ1 cases against witness graphs generated by VCF-RDFizer v3.0.3~\cite{vcfrConverter2026}, keeping fixtures, queries, and expected answers unchanged.
This converter executes RML mappings~\cite{dimou2014rml} with RMLStreamer~\cite{haesendonck2019rmlstreamer}, whereas the VCF Core materializer constructs triples via a Python script.
The programs share an author, so agreement cannot rule out a shared interpretation error.
The comparison can nevertheless reveal implementation-specific omissions or incorrect associations, and the round-trip separately checks recovered values against the source VCF.
Every fixture is converted in both profiles, and both producers' graphs are generated during the run.
For each case, axis, and profile, we record whether both pass, either producer alone passes, or both fail.
We also compare the two profiles case by case.
Because producers mint different resource bases, fragment-bearing IRIs are compared by fragment alone; a wrong namespace with the correct fragment could therefore escape detection.

\paragraph{Semantic round-trip.}
A complementary check, run on expanded graphs only, reconstructs the seven fixed columns, sample genotypes, header lines, INFO entries, and non-GT FORMAT values from structured RDF properties and compares them with the source.
The preserved-source properties \term{infoRaw}, \term{sampleDataRaw}, and \term{genotypeString} are not used.
Deliberate alterations or deletions of genotype, field, and header data verify that mismatches are detected.
A leading phase indicator is normalized because VCF 4.4--4.5 permit equivalent forms such as \texttt{|0|1} and \texttt{0|1}, but allele order, missing calls, ploidy, and intervening phase indicators must still match exactly.
A condensed round-trip was not run: the round-trip check has no decoder for the encoded sample vectors yet, although RQ1's preservation queries show that those values can be decoded.
Only sample-level values would differ, since both profiles agree on every query above the sample level (Section~\ref{sec:profilesresult}).

\subsection{Results}
\label{sec:crossproducer}

Re-executing the 210 cases on two axes and two profiles gives \textbf{840 checks}: \textbf{693 pass for both producers}, \textbf{45 fail for both}, and \textbf{102 differ} across 15 requirements.
Every disagreement is a materializer pass and converter failure; none shows the inverse.
The 45 shared failures are condensed-structure queries requiring per-sample resources.
Of the 53 exercised requirements, 38 agree on every check and 32 pass on every check.
Inspection of each disagreement identifies the causes in Table~\ref{tab:divergence}.
No observed divergence requires a representational mechanism absent from the pinned vocabulary; every missing property is already declared.

\begin{table}[htbp]
\caption{Causes of disagreement across 15 requirements.}
\label{tab:divergence}
\centering\small
\begin{tabularx}{\linewidth}{@{}>{\raggedright\arraybackslash}p{0.29\linewidth}c>{\raggedright\arraybackslash}X@{}}
\toprule
Cause & Reqs. & Finding \\
\midrule
Unemitted properties & 12 & Twenty declared properties absent, including accessors for gVCF blocks, breakend mates, and allele indices. \\
Incomplete decomposition & 1 & Tandem-repeat summary emitted without per-allele components. \\
Permitted datatype difference & 1 & QUAL uses \term{vcfc:VCFFloat} or \term{xsd:decimal}; both satisfy the shape but differ under \term{DATATYPE()}. \\
Encoding convention & 1 & Different handling of an empty local-allele list (R27). \\
\bottomrule
\end{tabularx}
\end{table}

An earlier comparison also exposed a silent local-allele error, missed by the converter's existing tests: the converter associated \term{Number=LA} and \term{Number=LR} values with record alleles positionally instead of resolving the sample's local subset.
Under \texttt{LAA=2,4}, depths \texttt{20,30,10} belong to global alleles 0, 2, and 4, but were attached to 0, 1, and 2.
That defect and eight further divergences were repaired~\cite{vcfrConverter2026}.
In the evaluated release, all local-allele requirements pass for both producers on expanded-profile checks.

\paragraph{Semantic round-trip and profile comparison.}
Across 38 fixtures, reconstruction from structured properties recovers the seven fixed columns of all \textbf{109 records}, all \textbf{151 sample genotypes}, all \textbf{431 header lines}, and \textbf{400 of 401} INFO entries and non-GT FORMAT subfields.
The exception is an empty \texttt{LAA} list: the materializer writes an empty value, while the converter emits the field without a value triple.
This loses the distinction between a present, empty list, meaning reference-only local alleles, and a field without a value.
The round-trip thus provides a complementary check on R27.
The associated \texttt{LAD} item retains the \term{forAllele} link, so the error does not remove that depth association.

\label{sec:profilesresult}
The sample profiles give the same outcome for both producers on \textbf{352 of 420} case-and-axis comparisons.
All 68 disagreements concern FORMAT values or genotypes within samples, spanning 15 requirements.
Questions about files, headers, records, and alleles agree throughout.
The profiles therefore agree above the sample level; below it, condensed values must be decoded first.

